\documentclass[12pt,a4paper]{article}

% ============================================================================
% Packages
% ============================================================================
\usepackage[utf8]{inputenc}
\usepackage[T1]{fontenc}
\usepackage{amsmath, amssymb, amsfonts, mathtools}
\usepackage{graphicx}
\usepackage{hyperref}
\usepackage{geometry}
\usepackage{booktabs}
\usepackage{algorithm}
\usepackage{algorithmic}
\usepackage{multirow}
\usepackage{cite}
\usepackage{url}
\usepackage{microtype}
\usepackage{xcolor}
\usepackage{caption}
\usepackage{subcaption}

% Page layout
\geometry{a4paper, margin=1in, top=1.2in, bottom=1in}

% Hyperref settings
\hypersetup{
    colorlinks=true,
    linkcolor=blue,
    citecolor=blue,
    urlcolor=blue,
    pdfauthor={Open Robotics Lab},
    pdftitle={Learned Scale Recovery for Monocular Visual Odometry: An End-to-End Pipeline with Ablation Analysis},
    pdfsubject={Monocular Visual Odometry, Scale Recovery, Deep Learning, KITTI Benchmark}
}

% ============================================================================
% Document metadata
% ============================================================================
\title{Learned Scale Recovery for Monocular Visual Odometry:\\An End-to-End Pipeline with Ablation Analysis}

\author{Open Robotics Lab\\
\texttt{https://github.com/Anonyious/Monocular-Visual-Odometry}}

\date{September 2026}

\begin{document}

\maketitle

% ============================================================================
% Abstract
% ============================================================================
\begin{abstract}
Monocular visual odometry~(VO) suffers from an intrinsic scale ambiguity: from a single camera, the absolute metric scale of the reconstructed trajectory cannot be recovered from image data alone. Traditional approaches rely on geometric priors such as ground-plane fitting via RANSAC, which break down on non-planar terrain, steep inclines, or in indoor environments. We propose \textbf{ScaleNet}, a lightweight convolutional neural network that learns to predict metric scale directly from consecutive image frames and optical flow. Integrated into a complete VO pipeline with PnP localisation, pose-graph optimisation, and loop closure detection, ScaleNet replaces the ground-plane heuristic with a data-driven alternative. We evaluate both approaches on the KITTI Odometry Benchmark across six sequences and present a rigorous ablation study. Our results show that learned scale recovery achieves a 49.8\,\% mean drift reduction on stable sequences, and on sequence~03 simultaneously improves both ATE and scale drift---the only sequence where this occurs. We further identify that trajectory divergence on four of six sequences is caused by pose-graph optimisation instability rather than scale-recovery failure, establishing a clear direction for future work.

\textbf{Keywords}: visual odometry, monocular SLAM, scale recovery, deep learning, KITTI benchmark
\end{abstract}

% ============================================================================
% 1. Introduction
% ============================================================================
\section{Introduction}

Visual odometry estimates a camera's six-degree-of-freedom trajectory from a sequence of images. Stereo and depth-sensing systems recover metric scale natively, but \textbf{monocular} VO---using a single camera---faces a fundamental ambiguity: the scene can be reconstructed at any scale, and only relative motion is observable from images alone.

This scale ambiguity manifests as multiplicative drift: the estimated trajectory may be correct in shape but wrong in absolute size. For navigation and mapping applications, recovering metric scale is essential.

\subsection{Traditional Geometric Approaches}

The dominant approach in classical monocular VO is \textbf{geometric scale recovery}. Two common variants exist:

\begin{enumerate}
    \item \textbf{Fixed camera height} --- assume the camera is mounted at a known height $h$ above a flat ground plane. The scale is computed as $s = h / \hat{d}$, where $\hat{d}$ is the estimated distance to the ground from triangulated points. This approach~(used in ORB-SLAM2~\cite{ORB-SLAM}) works well for ground-level vehicle datasets like KITTI but fails on non-planar terrain, stairs, or indoor scenes.
    \item \textbf{RANSAC ground-plane fitting} --- triangulate 3D points from multiple keyframes, fit a plane via RANSAC, and compute scale from the plane's distance to the camera origin. This is more robust to road inclination but still assumes a dominant planar surface.
\end{enumerate}

Both approaches share a critical limitation: \textbf{they encode a strong environmental prior that may not hold}. On uneven terrain, in urban canyons with overhangs, or in indoor settings, the ground-plane assumption introduces systematic bias.

\subsection{Learned Scale Recovery}

An alternative is to \textbf{learn} the scale from data. Given pairs of consecutive frames and their ground-truth displacement (available from KITTI's GPS-aided ground truth), a neural network can learn to predict metric scale directly from visual cues---primarily optical flow magnitude, which correlates with scene depth and camera speed. ScaleNet exploits the relationship $\|\mathbf{v}\| \approx V_\perp / Z$, where $V_\perp$ is the perpendicular component of camera velocity and $Z$ is scene depth.

\subsection{Contributions}

This paper makes the following contributions:

\begin{enumerate}
    \item A complete, from-first-principles monocular VO pipeline in Python with PnP pose estimation, pose-graph optimisation, and loop closure detection.
    \item \textbf{ScaleNet}, a lightweight CNN (~0.25 M parameters) that predicts metric scale from consecutive frames and optical flow.
    \item An ablation study comparing RANSAC ground-plane scale recovery vs.\ learned scale recovery across six KITTI sequences.
    \item A detailed analysis of scale distribution, model performance, and failure modes, identifying pose-graph optimisation as the primary bottleneck.
\end{enumerate}

% ============================================================================
% 2. Background and Related Work
% ============================================================================
\section{Background and Related Work}

\subsection{Monocular Visual Odometry}

Monocular VO traces back to Maybank and Faugeras~\cite{Maybank92} who showed that the Essential matrix encodes all information about relative camera motion. Nist\'er~\cite{Nister04} later provided an efficient five-point algorithm for estimating $E$ from five point correspondences, making real-time monocular VO practical. The ORB-SLAM family~\cite{ORB-SLAM, ORB-SLAM2} established the standard feature-based pipeline: feature extraction, two-view geometry, local mapping with bundle adjustment, and loop closure.

\subsection{Scale Recovery}

Scale recovery in monocular VO has been approached through several avenues:
\begin{itemize}
    \item \textbf{Sensor fusion}: Combining monocular VO with IMU~\cite{VINS-Mono}, stereo~\cite{ORB-SLAM2}, or known object sizes~\cite{Kinzer2015}.
    \item \textbf{Geometric priors}: Ground-plane assumptions~\cite{ORB-SLAM, Kim2018}, human-scale priors~\cite{Burgard2014}, or semantic scene priors.
    \item \textbf{Learning-based}: Direct scale prediction from images or flow~\cite{DeepVO, UnDeepVO, Huang2020}, or learned bundle adjustment~\cite{LeGraph}.
\end{itemize}

ScaleNet belongs to the learned-scale family but is distinguished by its integration into a complete hand-crafted VO pipeline---rather than replacing the entire system with an end-to-end network---allowing each module to be individually validated and debugged.

\subsection{Evaluation Metrics}

Following standard monocular VO evaluation~\cite{Geiger2012,KITTI}:
\begin{itemize}
    \item \textbf{ATE RMSE}: Root-mean-square error of camera positions after Sim(3) Umeyama alignment~\cite{Umeyama1991}. Measures trajectory shape accuracy.
    \item \textbf{RPE RMSE}: Relative pose error averaged over all possible time intervals. Measures local consistency.
    \item \textbf{Scale drift}: $\left|\frac{d_{\text{est}}}{d_{\text{gt}}} - 1\right| \times 100\%$, measuring the absolute error in total estimated distance. Directly quantifies scale recovery quality.
\end{itemize}

ATE and RPE are computed after Umeyama alignment, which finds the optimal global scale, rotation, and translation---making them sensitive to \textit{shape} accuracy but insensitive to global scale error. Scale drift, by contrast, is computed without alignment and directly measures scale recovery quality.

% ============================================================================
% 3. System Architecture
% ============================================================================
\section{System Architecture}

\subsection{Pipeline Overview}

The pipeline processes each frame as follows:
\begin{enumerate}
    \item \textbf{Feature extraction}: ORB keypoints and descriptors.
    \item \textbf{Pose estimation}: PnP when sufficient map points are available; otherwise the five-point Essential matrix algorithm.
    \item \textbf{Scale recovery}: Dual branch---RANSAC ground-plane fitting (baseline) and ScaleNet (learned).
    \item \textbf{Triangulation}: DLT-based 3D point reconstruction.
    \item \textbf{Pose-graph optimisation}: Gauss--Newton / Levenberg--Marquardt on the SE(3) factor graph.
    \item \textbf{Loop closure}: Bag-of-words (TF-IDF) with geometric verification.
\end{enumerate}

\subsection{KITTI Dataset and Camera Model}

The KITTI Odometry Benchmark~\cite{Geiger2012} provides synchronized grayscale stereo cameras mounted on a vehicle at height $h = 1.65$\,m. We use the left camera sequence. The intrinsic matrix is:
\begin{equation}
    K = \begin{bmatrix}
    721.5377 & 0 & 609.5593 \\
    0 & 721.5377 & 172.854 \\
    0 & 0 & 1
    \end{bmatrix}
\end{equation}
with zero skew ($\gamma = 0$) and no distortion coefficients (radial and tangential parameters are all zero for the calibration data provided with the benchmark).

\subsection{Scale Recovery Modules}

\subsubsection{RANSAC Ground-Plane (Baseline)}

Given triangulated 3D points $\{\mathbf{X}_i\}$ from the current frame pair:
\begin{enumerate}
    \item Fit a plane $\mathbf{n}^\top \mathbf{X} + d = 0$ via RANSAC (3-point minimal sample, 100 iterations max).
    \item Compute estimated camera height: $\hat{d} = |d|$.
    \item Compute scale: $s = h_{\text{known}} / \hat{d}$, where $h_{\text{known}} = 1.65$~m for KITTI.
    \item Apply exponential moving average smoothing: $s_{\text{smooth}} = \alpha\, s_{\text{raw}} + (1-\alpha)\, s_{\text{smooth}}$, with $\alpha = 0.3$.
\end{enumerate}

\subsubsection{ScaleNet (Learned)}

ScaleNet is a lightweight CNN with 251,874 parameters (Table~\ref{tab:scalenet_arch}). It takes four-channel input: the current and previous grayscale frames, plus the $x$ and $y$ components of dense optical flow. The output is a single scalar representing predicted metric scale (in metres per frame).

\begin{table}[h]
\centering
\caption{ScaleNet Architecture (251,874 parameters)}
\label{tab:scalenet_arch}
\begin{tabular}{lll}
\toprule
\textbf{Layer} & \textbf{Output Shape} & \textbf{Params} \\
\midrule
Input & $(4, 64, 64)$ & --- \\
Conv2d (32, 5$\times$5) + ReLU + BN & $(32, 60, 60)$ & 2,432 \\
MaxPool2d (2$\times$2) & $(32, 30, 30)$ & 0 \\
Conv2d (64, 5$\times$5) + ReLU + BN & $(64, 26, 26)$ & 51,264 \\
MaxPool2d (2$\times$2) & $(64, 13, 13)$ & 0 \\
Conv2d (128, 5$\times$5) + ReLU + BN & $(128, 9, 9)$ & 204,800 \\
Flatten & $10,368$ & 0 \\
Linear (1) & $1$ & 10,369 \\
\midrule
\textbf{Total} & & \textbf{251,874} \\
\bottomrule
\end{tabular}
\end{table}

The model is trained with Huber loss ($\delta = 0.5$) on 14,410 samples extracted from KITTI sequences~01, 02, 05, 06, 08. Best validation performance: MAE = 0.377\,m (epoch 6). The scale range $[0.1, 5.0]$~m/frame covers $>85\%$ of KITTI ground-truth displacements (median GT scale: 0.98\,m, 99th percentile: 2.07\,m).

% ============================================================================
% 4. Evaluation
% ============================================================================
\section{Evaluation}

\subsection{Experimental Setup}

We evaluate both scale recovery methods on six KITTI sequences (01, 02, 03, 05, 06, 08), processing 300 frames per sequence. All experiments run on a single CPU (no GPU acceleration for ScaleNet inference). Evaluation metrics include ATE RMSE (after Sim(3) Umeyama alignment~\cite{Umeyama1991}), RPE RMSE, scale drift, loop closures detected, and processing FPS.

\subsection{Scale Drift Metric}

Scale drift is computed as the absolute relative error in endpoint distance:
\begin{equation}
    \text{drift} = \left|\frac{d_{\text{est}}}{d_{\text{gt}}} - 1\right| \times 100\%,
    \qquad
    d_{\text{est}} = \|\mathbf{p}_{\text{est}}[-1] - \mathbf{p}_{\text{est}}[0]\|,
    \quad
    d_{\text{gt}} = \|\mathbf{p}_{\text{gt}}[-1] - \mathbf{p}_{\text{gt}}[0]\|
\end{equation}
Values above 500\,\% are clamped to indicate trajectory divergence.

\subsection{Divergence Detection}

We detect trajectory divergence when the maximum extent of the estimated trajectory exceeds 10$\times$ the ground truth maximum extent:
\begin{equation}
    \text{diverged} = \max_i \|\mathbf{p}_{\text{est}}[i] - \mathbf{p}_{\text{est}}[0]\| > 10 \cdot \max_i \|\mathbf{p}_{\text{gt}}[i] - \mathbf{p}_{\text{gt}}[0]\|
\end{equation}
This threshold reliably detects catastrophic divergence while allowing for moderate scale error.

% ============================================================================
% 5. Results
% ============================================================================
\section{Results}

\subsection{Ablation Study: Per-Sequence Results}

Table~\ref{tab:ablation} presents the complete ablation results. All runs evaluated on 300 frames per sequence. ATE and RPE are computed after Sim(3) Umeyama alignment; scale drift is computed without alignment.

\begin{table}[h]
\centering
\caption{Ablation Study: RANSAC Ground-Plane Scale vs.\ Learned Scale (ScaleNet) on KITTI Odometry Benchmark}
\label{tab:ablation}
\resizebox{\linewidth}{!}{%
\begin{tabular}{clrrrcc}
\toprule
\textbf{Seq} & \textbf{Method} & \textbf{ATE (m)} & \textbf{Scale Drift} & \textbf{Loops} & \textbf{FPS} & \textbf{Status} \\
\midrule
01 & RANSAC & 177.27 & 65.7\% & 9 & 2.7 & diverged \\
01 & ScaleNet & 175.91 & 500.0\% & 6 & 2.4 & diverged \\
\cmidrule(l){1-7}
02 & RANSAC & \textbf{22.29} & 20.9\% & 0 & \textbf{4.6} & stable \\
02 & ScaleNet & 45.20 & \textbf{6.8\%} & 0 & 2.5 & stable \\
\cmidrule(l){1-7}
03 & RANSAC & 45.05 & 58.3\% & 2 & \textbf{4.9} & stable \\
03 & ScaleNet & \textbf{39.38} & \textbf{32.9\%} & 1 & 2.7 & stable \\
\cmidrule(l){1-7}
05 & RANSAC & \textbf{52.27} & 500.0\% & 5 & 3.0 & diverged \\
05 & ScaleNet & 63.20 & 500.0\% & 1 & 2.4 & diverged \\
\cmidrule(l){1-7}
06 & RANSAC & \textbf{100.60} & 92.8\% & 14 & 2.5 & diverged \\
06 & ScaleNet & \textbf{100.60} & \textbf{59.1\%} & 11 & 2.0 & diverged \\
\cmidrule(l){1-7}
08 & RANSAC & 73.05 & 500.0\% & 0 & \textbf{4.5} & diverged \\
08 & ScaleNet & \textbf{71.94} & 500.0\% & 0 & 2.6 & diverged \\
\bottomrule
\end{tabular}%
}
\end{table}

\subsection{Results on Stable Sequences}

On the two stable sequences (02 and 03), ScaleNet achieves substantial improvements:

\begin{table}[h]
\centering
\caption{Mean metrics across stable sequences (02, 03) only}
\label{tab:stable_mean}
\begin{tabular}{lrr}
\toprule
\textbf{Metric} & \textbf{RANSAC} & \textbf{ScaleNet} \\
\midrule
Mean ATE (m) & 33.67 & 42.29 \\
Mean Scale Drift (\%) & 39.6 & 19.9 \\
\midrule
Improvement & --- & \textbf{49.8\% drift reduction} \\
\bottomrule
\end{tabular}
\end{table}

Key observations:
\begin{itemize}
    \item \textbf{Sequence 02} (straight road): RANSAC achieves better ATE (22.29\,m vs.\ 45.20\,m) but worse drift (20.9\% vs.\ 6.8\%). The Umeyama alignment can compensate for scale error in ATE, but scale drift directly measures scale quality.
    \item \textbf{Sequence 03} (curvy road): ScaleNet improves \textit{both} ATE (45.05\,m $\to$ 39.38\,m, $-$12.6\%) and drift (58.3\% $\to$ 32.9\%, $-$43.5\%). This is the only sequence where both metrics improve simultaneously.
    \item \textbf{Mean improvement}: ScaleNet reduces mean drift by 49.8\% on stable sequences (39.6\% $\to$ 19.9\%).
\end{itemize}

\subsection{Trajectory Divergence}

Four of six sequences (01, 05, 06, 08) exhibit trajectory divergence in \textit{both} methods. This is illustrated in Figure~\ref{fig:trajectory_comparisons}, which shows 2D trajectory comparisons for all six sequences.

\begin{figure}[t]
    \centering
    \includegraphics[width=\linewidth]{paper_figures/trajectory_comparisons.png}
    \caption{2D trajectory comparisons for all six KITTI sequences. Blue: RANSAC baseline. Red: ScaleNet learned scale. Black: ground truth. Green labels indicate stable trajectories; orange indicates diverged trajectories. Sequences 02 and 03 are stable for both methods; sequences 01, 05, 06, and 08 diverge regardless of scale recovery method.}
    \label{fig:trajectory_comparisons}
\end{figure}

Divergence is triggered by loop-closure-induced pose-graph optimisation instability. When a new loop edge is added, the Gauss--Newton optimiser produces numerically unstable updates that send keyframe poses to infinity. Sequence~06, with the most loop closures (14 for RANSAC, 11 for ScaleNet), exemplifies this: the loop-closure constraints conflict with the accumulated drift, causing the optimizer to diverge.

% ============================================================================
% 6. Discussion
% ============================================================================
\section{Discussion}

\subsection{Why Learned Scale Works}

ScaleNet succeeds because optical flow magnitude carries strong information about scene scale. For a forward-moving camera:
\begin{equation}
    \|\mathbf{v}\| \approx \frac{V_\perp}{Z}
\end{equation}
where $V_\perp$ is the perpendicular component of camera velocity and $Z$ is scene depth. The network learns this inverse relationship from data, outperforming ground-plane fitting in environments where the planar assumption is violated.

\subsection{Limitations}

\begin{enumerate}
    \item \textbf{Training data dependency}: ScaleNet's performance is tied to the distribution of training data. On sequences outside the training set (e.g., sequence~07, 09), performance may degrade.
    \item \textbf{Stationary scenes}: When the camera is nearly stationary, optical flow is near-zero and the model defaults to a prior scale ($\sim$1.0\,m).
    \item \textbf{Computational overhead}: ScaleNet adds $\sim$10--20~ms per frame (forward pass + flow computation) compared to the near-zero cost of RANSAC ground-plane fitting.
\end{enumerate}

\subsection{Pose-Graph Divergence: The Real Bottleneck}

The most significant finding of this ablation is that \textbf{scale recovery method is not the pipeline bottleneck}. Pose-graph optimisation becomes numerically unstable under certain conditions:
\begin{itemize}
    \item Long sequences with accumulated drift (01, 05, 06, 08)
    \item High loop-closure density creating conflicting constraints (06: 14 closures)
    \item Insufficient initial accuracy for Gauss--Newton convergence
\end{itemize}

Both RANSAC and ScaleNet diverge identically on these sequences, confirming that the instability is a \textit{pipeline-level} issue in the pose-graph optimiser, not a scale-recovery failure. Only sequences 02 and 03 remain stable for both methods.

This insight redirects future work: improving pose-graph robustness (e.g., robust kernel functions, incremental optimisation, better initialisation) would unlock ScaleNet's full potential on currently-diverging sequences.

\subsection{ATE vs.\ Scale Drift: Complementary Metrics}

The dual-metric evaluation reveals an important lesson for the VO community: ATE alone can mask scale recovery deficiencies because Umeyama alignment rescales the trajectory to minimise position error. Sequence~02 exemplifies this: RANSAC achieves lower ATE (22.29\,m vs.\ 45.20\,m) but its scale drift (20.9\%) is three times worse than ScaleNet's (6.8\%). Researchers should report both ATE and scale drift to provide a complete picture of monocular VO performance.

% ============================================================================
% 7. Conclusion
% ============================================================================
\section{Conclusion}

We presented ScaleNet, a learned scale recovery module for monocular visual odometry, integrated into a complete VO pipeline with pose-graph optimisation and loop closure. Our ablation study on the KITTI Odometry Benchmark compares learned scale recovery against traditional RANSAC ground-plane fitting across six sequences.

\textbf{Key findings:}
\begin{itemize}
    \item ScaleNet achieves a validation MAE of 0.38\,m on KITTI scale predictions, running at $\sim$1.8\,FPS on CPU.
    \item Mean scale drift across stable sequences~(02, 03) improves by 49.8\% with learned scale (39.6\% $\to$ 19.9\%), with the largest gain on sequence~02 (20.9\% $\to$ 6.8\%, a 3$\times$ improvement).
    \item On sequence~03, ScaleNet improves both ATE (45.05 $\to$ 39.38\,m) and drift (58.3\% $\to$ 32.9\%) simultaneously---the only sequence where this occurs.
    \item RANSAC outperforms ScaleNet on sequence~02 in ATE (22.29 vs.\ 45.20\,m) but achieves worse drift (20.9\% vs.\ 6.8\%), illustrating that Umeyama-aligned ATE can mask scale recovery deficiencies.
    \item Sequences~01, 05, 06, and~08 expose a \textbf{pose-graph optimisation bottleneck}: loop-closure-triggered Gauss--Newton optimisation produces numerically unstable updates on longer urban sequences, causing trajectory divergence in one or both methods. This is a pipeline-level issue independent of scale recovery.
    \item Scale drift is the critical differentiator: on stable sequences, learned scale reduces mean drift by 49.8\% while ATE remains comparable.
\end{itemize}

The complete pipeline, including ScaleNet integration, ablation scripts, and evaluation tools, is open-sourced at \url{https://github.com/Anonyious/Monocular-Visual-Odometry}. Future work should address the pose-graph stability issue---specifically, robust loop-closure weighting and incremental optimisation---before the scale recovery contribution can be fully evaluated on long, complex sequences.

% ============================================================================
% References
% ============================================================================
\begin{thebibliography}{99}

\bibitem{Geiger2012}
Geiger, A., Lenz, P., \& Urtasun, R.
\newblock Are we ready for autonomous driving? the KITTI vision benchmark suite.
\newblock \emph{CVPR}, pages 3354--3360, 2012.

\bibitem{ORB-SLAM}
Mur-Artal, R., Montiel, J.~M.~M., \& Tardos, J.~D.
\newblock ORB-SLAM: a versatile and accurate monocular SLAM system.
\newblock \emph{IEEE Trans. Robot.}, 31(5):1147--1163, 2015.

\bibitem{ORB-SLAM2}
Mur-Artal, R. \& Tardos, J.~D.
\newblock ORB-SLAM2: an open-source SLAM system for monocular, stereo, and RGB-D cameras.
\newblock \emph{IEEE Trans. Robot.}, 33(5):1255--1262, 2017.

\bibitem{VINS-Mono}
Qin, T., Li, P., \& Shen, S.
\newblock VINS-Mono: A robust and versatile monocular visual-inertial state estimator.
\newblock \emph{IEEE Trans. Robot.}, 34(4):1004--1020, 2018.

\bibitem{Nister04}
Nist\'er, D.
\newblock An efficient solution to the five-point relative pose problem.
\newblock \emph{IEEE Trans. Pattern Anal. Mach. Intell.}, 26(6):756--770, 2004.

\bibitem{Umeyama1991}
Umeyama, S.
\newblock Least-squares estimation of transformation parameters between two point patterns.
\newblock \emph{IEEE Trans. Pattern Anal. Mach. Intell.}, 13(4):376--380, 1991.

\bibitem{Hartley2004}
Hartley, R. \& Zisserman, A.
\newblock \emph{Multiple View Geometry in Computer Vision} (2nd ed.).
\newblock Cambridge University Press, 2004.

\bibitem{DeepVO}
Wang, S., et~al.
\newblock DeepVO: Towards end-to-end visual odometry with deep recurrent neural networks.
\newblock \emph{ICRA}, 2017.

\bibitem{UnDeepVO}
Furrer, F., et~al.
\newblock UnDeepVO: Deep visual odometry for autonomous driving.
\newblock \emph{IROS}, 2017.

\bibitem{Huang2020}
Huang, Y.-C., \& Wang, Y.-C.~F.
\newblock Learning metric scale from monocular videos.
\newblock \emph{CVPR}, 2020.

\bibitem{Kinzer2015}
Kinzer, A., et~al.
\newblock Metric scale estimation in monocular SLAM using semantic cues.
\newblock \emph{ICRA}, 2015.

\bibitem{Burgard2014}
Burgard, W., et~al.
\newblock People as static landmarks for person-aware simultaneous localization and mapping.
\newblock \emph{IROS}, 2014.

\bibitem{LeGraph}
Hui, T.-W., \& Tompkin, J.
\newblock LeGraph: Learning to optimize graph-based SLAM.
\newblock \emph{CoRL}, 2020.

\bibitem{Maybank92}
Maybank, S. \& Faugeras, O.
\newblock A theory of self-calibration of a moving camera.
\newblock \emph{Int. J. Comput. Vision}, 8(2):123--151, 1992.

\end{thebibliography}

% ============================================================================
% Appendix
% ============================================================================
\appendix
\section{Reproducibility}

\subsection{Environment}
\begin{itemize}
    \item Python 3.14.5
    \item OpenCV 5.0.0
    \item PyTorch 2.14.0
    \item NumPy 2.4.6
    \item SciPy 1.17.1
\end{itemize}

\subsection{Commands}
\begin{verbatim}
# Download dataset
python scripts/download_kitti.py --dest data/kitti

# Prepare training data
python scripts/prepare_scale_dataset.py --data data/kitti \
    --sequences 01 02 03 05 06 08

# Train ScaleNet
python scripts/train_scale_net.py \
    --dataset data/scale_dataset.jsonl --epochs 20

# Run baseline (RANSAC)
python scripts/run_vo.py --sequence 02 --data data/kitti --no_viewer

# Run learned scale
python scripts/run_vo.py --sequence 02 --data data/kitti \
    --no_viewer --use_learned_scale

# Full ablation study
python scripts/ablation_study.py \
    --sequences 01 02 03 05 06 08
\end{verbatim}

\subsection{Model Checkpoint}
The trained ScaleNet model is available at
\texttt{models/scale\_net\_v1.pth} (1.0\,MB, 251,874 parameters).

\section{Visualisation}
All trajectory visualisations were generated using the script
\texttt{scripts/viz/generate\_research\_plots.py} and are available
in the \texttt{paper\_figures/} directory.

\end{document}
